%=====================================================================
\section*{Lecture 0: The map of the course}
\addcontentsline{toc}{section}{Lecture 0: The map of the course}
\label{sec:overview}

\paragraph{The problem.}
A small quantum system $S$ (here: a qubit) is coupled to a large environment $B$ (here: a bath of harmonic oscillators).
The full dynamics is unitary, but we only ever look at the system, so what we want is the \emph{reduced} dynamics: either the reduced density matrix $\rhoS(t)=\Tr_B\rho(t)$, or the expectation values and correlation functions of system operators.
Three formalisms were invented, by different communities and for different purposes, to do exactly this:
\begin{itemize}[itemsep=1pt]
\item \textbf{Keldysh} (Schwinger~\cite{Schwinger1961}, Keldysh~\cite{Keldysh1964}): write expectation values as a path integral or a diagrammatic series on a closed time contour that runs forward with $U$ and backward with $U^\dagger$.
  Natural for many-body systems, Green's functions and field theory.
\item \textbf{Nakajima--Zwanzig} (Nakajima~\cite{Nakajima1958}, Zwanzig~\cite{Zwanzig1960}) and its time-convolutionless (TCL) variant~\cite{Shibata1977,Chaturvedi1979}: project the full density matrix onto the ``relevant'' part $\rhoS\otimes\rhoB$ and derive a closed, exact integro-differential equation for $\rhoS(t)$.
  Natural for small systems and for questions about positivity, memory and master equations.
\item \textbf{Mori} (Mori~\cite{Mori1965}, Zwanzig~\cite[Ch.~8]{ZwanzigBook}): project \emph{operators} (Heisenberg picture) onto a chosen set of slow variables and derive a generalized Langevin equation.
  Natural for hydrodynamics, response functions and continued-fraction approximations.
\end{itemize}

\paragraph{The thesis.}
The three formalisms are not three theories; they are three bookkeeping schemes for the same perturbation series.
The precise statements, proved in the course, are:
\begin{keybox}[The three bridges]
\begin{enumerate}[itemsep=1pt]
\item Ordering superoperators along the real time axis \emph{is} Keldysh contour ordering (\cref{sec:lec3}).
  The ``classical/quantum'' Keldysh rotation is the pair of superoperators $A^\cl=\tfrac12\{A,\cdot\}$, $A^\qu=[A,\cdot]$.
\item The Nakajima--Zwanzig (NZ) memory kernel is the sum of \emph{irreducible} real-time diagrams of the reduced propagator, and the NZ equation is its Dyson equation (\cref{sec:lec5}).
  TCL is the same Dyson equation solved for a time-local generator (\cref{sec:lec6}).
\item Mori's memory function is the \emph{retarded self-energy} of the projected operator; for the full operator basis with the Heisenberg-adjoint inner product it is the adjoint of the NZ kernel (\cref{sec:lec8}).
\end{enumerate}
In all three, a Gaussian bath enters only through two kernels: the symmetrized noise $S(t)$ (Keldysh component $D^K$) and the response $\chi(t)$ (retarded component $D^R$).
\end{keybox}

\paragraph{The running example.}
Everything is worked out for
\begin{equation}
  H=\frac{\omega_0}{2}\sz+\sum_k\omega_kb_k^\dagger b_k
  +\big(\cos\theta\,\sz+\sin\theta\,\sx\big)\otimes\sum_kg_k\big(b_k+b_k^\dagger\big),
  \label{eq:model-intro}
\end{equation}
the spin-boson model with a general linear coupling.
The mixing angle $\theta$ lets us switch between the two textbook limits without changing the formalism: $\theta=0$ is pure dephasing (exactly solvable, the sensor's workhorse) and $\theta=\pi/2$ is transverse coupling (relaxation, the damped qubit).
Intermediate $\theta$ is unitarily equivalent to the biased dissipative two-level system of Leggett \emph{et al.}~\cite{Leggett1987}, as we show in \cref{sec:lec2}.
A Lorentzian spectral density makes the model numerically exact through the pseudomode mapping~\cite{Garraway1997,Tamascelli2018}, so each approximation in the course can be compared with the truth (\cref{sec:numerics}).

\paragraph{The route.}
\Cref{sec:lec1,sec:lec2} set up the language (Liouville space; the model and its bath).
\Cref{sec:lec3,sec:lec4} derive the Keldysh structure and the diagram rules from the von Neumann equation.
\Cref{sec:lec5,sec:lec6} derive NZ and TCL and identify them with the Dyson equation.
\Cref{sec:lec7} goes to the path integral: Feynman--Vernon, the polaron frame and the noninteracting-blip approximation.
\Cref{sec:lec8} does Mori and connects it to the other two.
\Cref{sec:lec9} rewrites the qubit problem as a Keldysh field theory, so the reader can recognize the structures in the many-body literature.
\Cref{sec:lec10} turns the machinery to quantum sensing: which protocol measures which bath kernel, and why non-Markovian memory changes precision scaling.

\paragraph{Prerequisites.}
Quantum mechanics at the level of second quantization for bosons, the interaction picture and the Dyson series; a first exposure to density matrices.
No prior knowledge of Keldysh, projection operators or path integrals is assumed.
Units: $\hbar=k_B=1$.

\clearpage
